\subsection{Local Adaptation of Global Watermarking Methods}
\label{app:local-adaptations}

This section describes how we adapt three global watermarking methods, OmniGuard, TrustMark, and PixelSeal, to the local watermarking setting evaluated in this work.
For each method, we explain how we restrict payload embedding to a selected image region and how the payload is subsequently recovered during decoding.
We use the released pretrained weights without additional training or fine-tuning.
During embedding, a binary mask specifies the selected region.
We refer to the pixelwise difference between an image before and after embedding as the \emph{watermark residual}: it records the pixel changes introduced by the embedder.

\paragraph{OmniGuard.}
OmniGuard embeds two distinct signals: an auxiliary signal for tamper localization and a binary payload for copyright identification.
We first apply its auxiliary embedding stage across the entire image.
We then resize this intermediate image to \(256\times256\) and pass it through the binary payload encoder.
Subtracting the encoder's input from its output gives the changes introduced specifically by payload embedding.
We resize this payload residual back to the benchmark image resolution of \(512\times512\), retain it only within the selected mask, and add it to the intermediate image.
Consequently, the auxiliary signal remains across the image, while the binary payload signal is confined to the selected region.
For payload recovery, we resize the complete received image to \(256\times256\) and apply the payload decoder.

\paragraph{TrustMark.}
TrustMark operates on rectangular regions.
We therefore replace the selected mask with its smallest enclosing axis-aligned bounding box, crop that region, and resize the crop to the encoder's native resolution.
We compute the watermark residual from the encoder's output and subtract its spatial mean separately for each color channel.
We then resize the residual back to the crop's original dimensions and add it to the corresponding region in the original image.
During decoding, TrustMark's region detector first predicts bounding boxes from the received image.
The predicted regions are cropped and resized to the decoder's native resolution for payload extraction.
We use the Q variant with a 100-bit payload and error correction disabled.

\paragraph{PixelSeal.}
We first apply PixelSeal's embedder to the complete image to obtain a globally watermarked output.
We subtract the original image from this output, multiply the resulting residual by the selected mask, and add the masked residual back to the original image.
This confines all changes introduced by PixelSeal to the selected region.
Unlike OmniGuard, this adaptation does not retain a separate auxiliary watermark outside that region.
We use the released 256-bit model and decode the payload from the complete received image using its whole-image decoder.

For all three methods, the selected mask is used only during embedding.
Decoding receives neither the ground-truth mask nor its bounding box.
